% ============================================================
%  PROBLEMAS DE PROBABILIDAD — UNIDAD 2
%  Un archivo por unidad; agregar bloques \dayheader según avanza el curso.
%  Compilar con pdflatex (funciona en Overleaf sin configuración extra).
% ============================================================

\documentclass[12pt, a4paper]{article}

% ---- codificación e idioma ---------------------------------
\usepackage[utf8]{inputenc}
\usepackage[T1]{fontenc}
\usepackage[spanish]{babel}

% ---- matemática --------------------------------------------
\usepackage{amsmath, amssymb, amsthm}

% ---- código R ----------------------------------------------
\usepackage{listings}
\usepackage{xcolor}

\definecolor{codegreen}{rgb}{0.2,0.5,0.2}
\definecolor{codegray}{rgb}{0.5,0.5,0.5}
\definecolor{codeblue}{rgb}{0.1,0.1,0.8}
\definecolor{backcolor}{rgb}{0.97,0.97,0.97}

\lstdefinestyle{Rstyle}{
  language=R,
  backgroundcolor=\color{backcolor},
  basicstyle=\ttfamily\small,
  keywordstyle=\color{codeblue}\bfseries,
  commentstyle=\color{codegreen}\itshape,
  stringstyle=\color{orange},
  numberstyle=\tiny\color{codegray},
  numbers=left,
  stepnumber=1,
  frame=single,
  rulecolor=\color{codegray},
  breaklines=true,
  showstringspaces=false,
  tabsize=2,
  captionpos=b
}
\lstset{style=Rstyle}

% ---- diseño de página --------------------------------------
\usepackage{geometry}
\geometry{margin=2.5cm}

\usepackage{enumitem}
\usepackage{parskip}
\usepackage{titlesec}
\usepackage{fancyhdr}
\usepackage{lastpage}
\usepackage{hyperref}


% ---- contador de problemas ---------------------------------
\newcounter{problema}
\renewcommand{\theproblema}{\arabic{diaactual}.\arabic{problema}}

\newcounter{diaactual}

\newcommand{\problema}[1][]{%
  \refstepcounter{problema}%
  \medskip
  \noindent\textbf{Problema~\theproblema%
    \ifx&#1&\else~—~#1\fi}%
  \par\smallskip
}

% ---- comando \dayheader ------------------------------------
\newcommand{\dayheader}[2]{%
  \setcounter{diaactual}{#1}%
  \setcounter{problema}{0}%
  \section*{Clase #1 \normalfont\small— #2}%
  \addcontentsline{toc}{section}{Clase #1: #2}%
}

% ---- nombre de la unidad (editar aquí) ---------------------
\newcommand{\nombreunidad}{Esperanza y Varianza}

% ============================================================
\begin{document}
% ============================================================

\textbf{8/9/26 - esperanza y varianza de variables aleatorias discretas}

\textbf{Problema 1}

En los videojuegos de Pokémon de primera generación, el  \textit{Ataque Furia}
golpea aleatoriamente al rival entre 2 y 5 veces, con las siguientes probabilidades:

\[
  P(G = 2) = P(G = 3) = \frac{3}{8}, \qquad P(G = 4) = P(G = 5) = \frac{1}{8}.
\]

Sea $G$ la variable aleatoria que cuenta el número de golpes del ataque.

\begin{enumerate}[label=\alph*)]
  \item Calcular $E(G)$ y $\mathrm{Var}(G)$.
  \item Cada golpe hace 15 puntos de daño con probabilidad $0{,}85$
        (y 0 puntos si no). Suponiendo $G = 5$, calcular el daño promedio
        esperado.
\end{enumerate}

\vspace{1 cm}

\textbf{Problema 2}

El número de pacientes que ingresa a una guardia de un centro de salud en una hora,
$X$, sigue una distribución de Poisson de parámetro $\lambda = 4$, es decir:
\[
  P(X = k) = \frac{e^{-4}\,4^k}{k!}, \quad k = 0, 1, 2, \ldots
\]

\begin{enumerate}[label=\alph*)]
  \item Calcular $P(X=0)$, $P(X=3)$ y $P(X \geq 2)$.
  \item ¿Cuál es el número esperado de pacientes por hora?
  \item Se dice que la guardia está \emph{saturada} si $X \geq 6$.
        Calcular la probabilidad de que la guardia esté saturada
        más de 5 horas en un día de 24 horas.
  \item Cada paciente paga \$3\,000 de coseguro. Si la guardia está saturada,
        los pacientes que llegan deciden irse para su casa y no pagan.
        ¿Cuál es el ingreso promedio esperado por hora?
\end{enumerate}

\vspace{1 cm}

\textbf{Problema 3}

Consideremos la variable aleatoria $X$ con la siguiente distribución de probabilidad:

\[
  \begin{array}{c|ccccc}
    x      & 0            & 1            & 2            & 3            & 4 \\
    \hline
    P(X=x) & \tfrac{1}{10} & \tfrac{2}{10} & \tfrac{3}{10} & \tfrac{3}{10} & \tfrac{1}{10}
  \end{array}
\]

\begin{enumerate}[label=\alph*)]
  \item Calcular $E(X)$ y $\mathrm{Var}(X)$.
  \item Sea $Y = 2X - 3$. Calcular $E(Y)$ y $\mathrm{Var}(Y)$
        usando las propiedades de esperanza y varianza.
  \item Sea $Z = (X-2)^2$. Calcular $E(Z)$.
\end{enumerate}

% ============================================================
%  Agregar más bloques \dayheader{2}{...} a medida que avanza la unidad.
% ============================================================

\end{document}